
Calibration -- the correspondence between a model's expressed confidence and its empirical accuracy -- is a central dimension of language model trustworthiness. Well-calibrated models provide uncertainty signals that can inform downstream decision-making: high confidence should correlate with high accuracy. Guo et al. (2017) established that modern neural networks tend to exhibit systematic overconfidence after training, and subsequent work has confirmed that RLHF alignment substantially amplifies this effect. Xie et al. (2024) reported that LLaMA-2-Chat achieves ECE values of 0.298 on MMLU and 0.507 on TruthfulQA, substantially exceeding its base model counterparts.

While the existence of alignment-induced miscalibration is well-documented, the mechanistic cause remains unresolved. Three distinct hypotheses have been proposed in the literature, each with different implications for corrective interventions:

\textbf{H1 (Scale Inflation).} Alignment amplifies confidence in the model's existing answer preferences without reordering them. Under H1, the Spearman rank correlation between base and aligned log-probability vectors should remain high (rho >= 0.90), and temperature scaling is expected to be an effective correction.

\textbf{H2 (Decision-Boundary Restructuring).} Alignment fundamentally changes which answer option the model selects, redistributing probability mass across options rather than merely inflating confidence margins. Under H2, Spearman rho should fall below 0.85, and temperature scaling may be insufficient because the underlying answer distribution has changed.

\textbf{H3 (Framing Susceptibility).} Alignment makes confidence allocation sensitive to question framing, producing larger miscalibration on tasks with adversarially designed alternatives. Under H3, calibration degradation on TruthfulQA (which contains plausible-but-false alternatives) should equal or exceed degradation on MMLU.

Prior work has not discriminated among these mechanisms. Xie et al. (2024) proposed Adaptive Temperature Scaling (ATS) as a correction but did not test whether the underlying mechanism is H1, H2, or H3. Li et al. (2024) used the Pythia alignment ladder to study RLHF effects on trustworthiness dimensions including toxicity and bias, but did not measure calibration. No study has conducted a pre-registered, multi-model mechanistic discrimination experiment for alignment-induced miscalibration.

This study addresses that gap using the Pythia model family (1.4B, 2.8B, 6.9B) with SFT, DPO, and PPO alignment variants -- a controlled setting where all variants share identical pretraining data and architecture. The contributions are as follows:

\begin{enumerate}
\item A mechanistic discrimination result: within the Pythia 1.4B-6.9B family, the data are consistent with H2 (decision-boundary restructuring) rather than H1 (scale inflation). All nine alignment-size pairs fall below the H1 threshold; eight of nine fall below the H2 diagnostic threshold.
\end{enumerate}

\begin{enumerate}
\item An exploratory ordering observation: DPO produces larger calibration degradation than PPO in all three Pythia sizes, contrary to the pre-registered prediction. This observation is based on public fallback checkpoints and should be treated as preliminary.
\end{enumerate}

\begin{enumerate}
\item H3 exclusion in the softmax-ECE setting: framing susceptibility is not supported as a primary driver, as MMLU calibration degradation exceeds TruthfulQA degradation for all alignment types.
\end{enumerate}

\begin{enumerate}
\item A reusable pre-registered methodology (Spearman rho thresholds, argmax partition, cross-benchmark diagnostic) applicable to any model family.
\end{enumerate}
